% Part: computability
% Chapter: computability-theory
% Section: halting-problem

\documentclass[../../../include/open-logic-section]{subfiles}

\begin{document}

\olfileid[tr]{cmp}{thy}{hlt}
\olsection{Durma Problemi}

Kuruluşu gereği evrensel kısmi hesaplanabilir $\fn{Un}(e,x)$ fonksiyonu,
ancak ve ancak $e$ ile kodlanan fonksiyonun hesabı $x$ girdisinde bir değer
üretirse tanımlıdır. Bunun böyle olup olmadığına karar verip veremeyeceğimizi
sormak doğaldır. Gerçekte veremeyiz. Turing makinesi hesaplama modelinde bu,
belirli bir Turing makinesinin belirli bir girdide durup durmadığının
hesaplama yoluyla kararlaştırılamadığı anlamına gelir. Bu nedenle aşağıdaki
teorem ``durma probleminin karar verilemezliği'' diye bilinir. Aşağıda iki
ispat vereceğiz. İlki önceki tartışmamızın çizgisini sürdürür, ikincisi daha
doğrudandır.

\begin{thm}
\ollabel{thm:halting-problem}
\[
  h(e,x)=\begin{cases}
    1 & \text{$\fn{Un}(e,x)$ tanımlıysa}\\
    0 & \text{aksi durumda}
  \end{cases}
\]
olsun. O zaman $h$ hesaplanabilir değildir.
\end{thm}

\begin{proof}
$h$ hesaplanabilir olsun. Bunun evrensel bir hesaplanabilir fonksiyon
tanımlamamızı sağlayacağını göstereceğiz. Şöyle tanımlayalım:
\[
  \fn{Un'}(e,x)=\begin{cases}
    \fn{Un}(e,x) & \text{$h(e,x)=1$ ise}\\
    0 & \text{aksi durumda.}
  \end{cases}
\]
Artık $\fn{Un'}(e,x)$ tümel bir fonksiyondur ve $h$ hesaplanabilirse o da
hesaplanabilirdir. Örneğin $g$ fonksiyonunu ilkel özyinelemeyle
\begin{align*}
  g(0,e,x) & \simeq 0\\
  g(y+1,e,x) & \simeq \fn{Un}(e,x)
\end{align*}
olarak, ardından
\[
  \fn{Un'}(e,x)\simeq g(h(e,x),e,x)
\]
biçiminde tanımlayabiliriz. $\fn{Un'}(e,x)$, $\fn{Un}(e,x)$ tanımlı olduğu
her yerde onunla aynı olduğundan $\fn{Un'}$, tümel olan kısmi hesaplanabilir
fonksiyonlar için evrenseldir. Bu, \olref[nou]{thm:no-univ} ile çelişir.
\end{proof}

\begin{proof}
$h(e,x)$ hesaplanabilir olsun. $g$ fonksiyonunu
\[
  g(x)=\begin{cases}
    0 & \text{$h(x,x)=0$ ise}\\
    \fundefined & \text{aksi durumda}
  \end{cases}
\]
olarak tanımlayalım. $g$ kısmi hesaplanabilirdir; örneğin
$\umin{y}{h(x,x)=0}$ olarak tanımlanabilir. Dolayısıyla bazı $e$ için her
$x$'te $g(x)\simeq\fn{Un}(e,x)$ olur. $g$, $e$'de tanımlı mıdır? Tanımlıysa
$g$ tanımı gereği $h(e,e)=0$ olur. $h$ tanımı gereği bu,
$\fn{Un}(e,e)$ ifadesinin tanımsız olduğu anlamına gelir. Her $x$ için
$g(x)\simeq\fn{Un}(e,x)$ varsayımımızdan $g(e)$'nin tanımsız olduğu çıkar;
bu bir çelişkidir. Öte yandan $g(e)$ tanımsızsa $h(e,e)\neq0$, dolayısıyla
$h(e,e)=1$ olur. Buradan $\fn{Un}(e,e)$ tanımlıdır. Fakat
$g(x)\simeq\fn{Un}(e,x)$ olduğundan $g(e)$ de tanımlı olurdu. Yine çelişki.
\end{proof}

\begin{tagblock}{TMs}
\begin{explain}
Bu savı Turing makineleri cinsinden anlatabiliriz. Bir Turing makinesi $E$'nin
betimi ile bir $x$ girdisini alıp $E$'nin $x$ girdisinde durup durmadığına
karar veren bir $H$ Turing makinesi bulunduğunu varsayalım. O zaman tek bir
$x$ girdisi alan, indisi $x$ olan $M_x$ makinesinin $x$ girdisinde durup
durmadığına karar vermek için $H$'yi çalıştıran ve tersini yapan başka bir
$G$ Turing makinesi kurabilirdik. Başka bir deyişle $H$, $M_x$ makinesinin
$x$ girdisinde durduğunu bildirirse $G$ sonsuz bir döngüye girer; $H$,
$M_x$ makinesinin $x$ girdisinde durmadığını bildirirse $G$ durur. $G$ kendi
indisini girdi olarak aldığında durur mu? Yukarıdaki sav, ancak ve ancak
durmazsa duracağını gösterir; bu bir çelişkidir. Dolayısıyla böyle bir
$H$ Turing makinesinin bulunduğu varsayımı yanlıştır.
\end{explain}
\end{tagblock}

\end{document}

